% ============================================================================
% Chapter 8 solutions: Simplex - Matrix Calculations
% Book source: part1-linear-programming/ch06-simplex/simplex-matrix-version.tex
% Exercise numbering follows \begin{ex} order in that file (9 exercises, all
% in the Exercises section; no embedded mid-chapter exercise exists in the
% current source).
% All basis computations verified with exact-fraction linear algebra and
% scipy.optimize.linprog, July 2026.
% ============================================================================
\section{Bab 8: Simpleks -- Perhitungan Matriks}

Bab 8 menyajikan kembali metode simpleks sebagai aljabar basis: $A_B$,
$A_B^{-1}\b$, dan biaya tereduksi
$\c_N^\top - \c_B^\top A_B^{-1} A_N$. Kesembilan latihannya terbagi menjadi
tiga pemanasan tentang penghitungan solusi basis dari suatu basis (8.1 sampai
8.3), tiga perhitungan inti dengan biaya tereduksi dan rumus kamus (8.4 sampai
8.6), dua latihan konseptual yang menghubungkan biaya tereduksi dengan konvensi
tanda pada kamus dan tableau (8.7, 8.8), serta satu pembuktian tantangan tentang
keunikan optimum (8.9). Setiap invers matriks, solusi basis, dan vektor biaya
tereduksi di bawah ini diverifikasi dengan perhitungan pecahan eksak. Solusi
pilihan dalam buku tersedia untuk 8.1, 8.4, dan 8.9 serta direproduksi di sini.

\exsol{8.1}{Matriks basis dan solusi basis}{1}

(Solusi pilihan dalam buku, telah diverifikasi.) Kolom $x_2$ dan $x_3$
memberikan
\[
A_B =
\begin{bmatrix}
2 & 1 \\
1 & 2
\end{bmatrix},
\qquad
A_B^{-1} = \frac{1}{3}
\begin{bmatrix}
2 & -1 \\
-1 & 2
\end{bmatrix}.
\]
Solusi basisnya adalah
\[
\x_B = A_B^{-1}\b = \frac{1}{3}
\begin{bmatrix}
2 & -1 \\
-1 & 2
\end{bmatrix}
\begin{bmatrix}
4 \\ 5
\end{bmatrix}
= \frac{1}{3}
\begin{bmatrix}
3 \\ 6
\end{bmatrix}
=
\begin{bmatrix}
1 \\ 2
\end{bmatrix},
\]
sehingga $(x_1, x_2, x_3) = (0, 1, 2)$. Kedua variabel basis tak negatif,
sehingga ini merupakan solusi basis layak.

\exsol{8.2}{Solusi basis yang tidak layak}{1}

Kolom $x_1$ dan $x_3$ memberikan
\[
A_B =
\begin{bmatrix}
1 & 1 \\
3 & 2
\end{bmatrix},
\qquad
\det A_B = 2 - 3 = -1,
\qquad
A_B^{-1} =
\begin{bmatrix}
-2 & 1 \\
3 & -1
\end{bmatrix}.
\]
Dengan menetapkan $x_2 = 0$, solusi basisnya adalah
\[
\x_B = A_B^{-1}\b =
\begin{bmatrix}
-2 & 1 \\
3 & -1
\end{bmatrix}
\begin{bmatrix}
4 \\ 5
\end{bmatrix}
=
\begin{bmatrix}
-3 \\ 7
\end{bmatrix},
\]
sehingga $(x_1, x_2, x_3) = (-3, 0, 7)$. Vektor ini memenuhi $A\x = \b$
(periksa:
$-3 + 2(0) + 7 = 4$ dan $3(-3) + 0 + 2(7) = 5$), tetapi variabel basis
$x_1 = -3$ melanggar kendala nonnegativitas $\x \geq \mathbf{0}$. Sebuah solusi
basis baru merupakan solusi basis \emph{layak} jika $\x_B \geq \mathbf{0}$,
sehingga basis ini memberikan solusi basis yang tak layak.

\exsol{8.3}{Basis keempat untuk contoh tiga basis}{1}

Kolom $x_2$ dan $s_2$ memberikan
\[
A_B =
\begin{bmatrix}
2 & 0 \\
3 & 1
\end{bmatrix},
\qquad
A_B^{-1} = \frac{1}{2}
\begin{bmatrix}
1 & 0 \\
-3 & 2
\end{bmatrix}
=
\begin{bmatrix}
\tfrac12 & 0 \\
-\tfrac32 & 1
\end{bmatrix}.
\]
Solusi basisnya adalah
\[
\x_B =
\begin{bmatrix} x_2 \\ s_2 \end{bmatrix}
= A_B^{-1}\b =
\begin{bmatrix}
\tfrac12 & 0 \\
-\tfrac32 & 1
\end{bmatrix}
\begin{bmatrix}
16 \\ 45
\end{bmatrix}
=
\begin{bmatrix}
8 \\ 21
\end{bmatrix}.
\]
Kedua komponennya tak negatif, sehingga basis tersebut \emph{layak}. Karena
$x_1 = s_1 = 0$, titik yang sesuai adalah $(x_1, x_2) = (0, 8)$:
titik sudut tempat batas $x_1 + 2x_2 = 16$ bertemu dengan sumbu $x_2$. (Basis
inilah yang ditunjukkan optimal pada contoh terurai dalam bagian optimalitas.)

\exsol{8.4}{Biaya tereduksi pada dua basis}{2}

(Solusi pilihan dalam buku, telah diverifikasi, dengan menampilkan
matriks-matriks antara.)

\textbf{Basis $B_1 = \{y, s_1, s_2\}$.} Dengan mengurutkan variabel basis
$(y, s_1, s_2)$,
\[
A_{B_1} =
\begin{bmatrix}
1 & 1 & 0 \\
1 & 0 & 1 \\
2 & 0 & 0
\end{bmatrix},
\qquad
A_{B_1}^{-1} =
\begin{bmatrix}
0 & 0 & \tfrac12 \\
1 & 0 & -\tfrac12 \\
0 & 1 & -\tfrac12
\end{bmatrix},
\qquad
\x_{B_1} = A_{B_1}^{-1}\b = \begin{bmatrix} 7 \\ 2 \\ 9 \end{bmatrix},
\]
sehingga $y = 7$, $s_1 = 2$, $s_2 = 9$: suatu solusi basis layak dengan nilai
objektif $\c_{B_1}^\top \x_{B_1} = 3 \cdot 7 = 21$. Dengan
$\c_{B_1}^\top = (3, 0, 0)$, pengali simpleksnya adalah
$\c_{B_1}^\top A_{B_1}^{-1} = (0, 0, \tfrac32)$, dan untuk variabel nonbasis
$x$ dan $s_3$ (kolom $(1,2,1)^\top$ dan $(0,0,1)^\top$),
\[
\c_N^\top - \c_{B_1}^\top A_{B_1}^{-1} A_N
= \begin{bmatrix} 2 & 0 \end{bmatrix}
- \begin{bmatrix} \tfrac32 & \tfrac32 \end{bmatrix}
= \begin{bmatrix} \tfrac12 & -\tfrac32 \end{bmatrix}.
\]
Biaya tereduksi untuk $x$ adalah $\tfrac12 > 0$, sehingga $B_1$
\emph{tidak} optimal; hasil ini sesuai dengan kamus kedua dari Bab~7,
$z = 21 + \tfrac12 x - \tfrac32 s_3$.

\textbf{Basis $B_2 = \{x, y, s_2\}$.} Dengan mengurutkan variabel basis
$(x, y, s_2)$,
\[
A_{B_2} =
\begin{bmatrix}
1 & 1 & 0 \\
2 & 1 & 1 \\
1 & 2 & 0
\end{bmatrix},
\qquad
A_{B_2}^{-1} =
\begin{bmatrix}
2 & 0 & -1 \\
-1 & 0 & 1 \\
-3 & 1 & 1
\end{bmatrix},
\qquad
\x_{B_2} = A_{B_2}^{-1}\b = \begin{bmatrix} 4 \\ 5 \\ 3 \end{bmatrix},
\]
suatu solusi basis layak dengan nilai $2 \cdot 4 + 3 \cdot 5 = 23$. Dengan
$\c_{B_2}^\top = (2, 3, 0)$, pengalinya adalah
$\c_{B_2}^\top A_{B_2}^{-1} = (1, 0, 1)$, dan untuk variabel nonbasis
$s_1, s_3$ (kolom $(1,0,0)^\top$ dan $(0,0,1)^\top$),
\[
\c_N^\top - \c_{B_2}^\top A_{B_2}^{-1} A_N
= \begin{bmatrix} 0 & 0 \end{bmatrix}
- \begin{bmatrix} 1 & 1 \end{bmatrix}
= \begin{bmatrix} -1 & -1 \end{bmatrix} \leq \mathbf{0}^\top,
\]
sehingga $B_2$ optimal, sesuai dengan kamus akhir
$z = 23 - s_1 - s_3$.

\exsol{8.5}{Basis mana yang layak?}{2}

\textbf{(a)} Keenam pasangan kolom dari
$A = \begin{bmatrix} 1 & 2 & 1 & 0 \\ 5 & 3 & 0 & 1 \end{bmatrix}$ memiliki
determinan
\[
\begin{aligned}
&\det A_{\{x_1,x_2\}} = -7,\qquad
\det A_{\{x_1,s_1\}} = -5,\qquad
\det A_{\{x_1,s_2\}} = 1,\\
&\det A_{\{x_2,s_1\}} = -3,\qquad
\det A_{\{x_2,s_2\}} = 2,\qquad
\det A_{\{s_1,s_2\}} = 1,
\end{aligned}
\]
semuanya tak nol, sehingga setiap pasangan kolom bebas linier dan setiap pilihan
merupakan suatu basis.

\textbf{(b)} Dengan menyelesaikan $A_B \x_B = \b$ untuk setiap basis
(variabel nonbasis sama dengan nol), diperoleh:
\begin{center}
\begin{tabular}{l|l|l|c}
Basis & Solusi basis & Titik $(x_1, x_2)$ & Layak? \\
\hline
$\{x_1, x_2\}$ & $x_1 = 6,\ x_2 = 5$ & $(6, 5)$ & ya \\
$\{x_1, s_1\}$ & $x_1 = 9,\ s_1 = 7$ & $(9, 0)$ & ya \\
$\{x_1, s_2\}$ & $x_1 = 16,\ s_2 = -35$ & $(16, 0)$ & tidak ($s_2 < 0$) \\
$\{x_2, s_1\}$ & $x_2 = 15,\ s_1 = -14$ & $(0, 15)$ & tidak ($s_1 < 0$) \\
$\{x_2, s_2\}$ & $x_2 = 8,\ s_2 = 21$ & $(0, 8)$ & ya \\
$\{s_1, s_2\}$ & $s_1 = 16,\ s_2 = 45$ & $(0, 0)$ & ya \\
\end{tabular}
\end{center}
Keempat basis layak bersesuaian dengan empat titik sudut daerah layak,
$(0,0)$, $(9,0)$, $(6,5)$, $(0,8)$; kedua basis tak layak merupakan titik
potong batas kendala yang berada di luar daerah tersebut.

\textbf{(c)} Ambil
\[
A = \begin{bmatrix} 1 & 2 & 1 & 0 \\ 2 & 4 & 0 & 1 \end{bmatrix}.
\]
Kedua kolom pertama memenuhi $(2, 4)^\top = 2 \cdot (1, 2)^\top$, sehingga
pasangan $\{1, 2\}$ bergantung linier ($\det = 0$) dan tidak membentuk basis.

\exsol{8.6}{Membangun kembali kamus dari rumus}{2}

Dengan urutan variabel $(x, y, s_1, s_2)$,
\[
A = \begin{bmatrix} 3 & 1 & 1 & 0 \\ 5 & 3 & 0 & 1 \end{bmatrix},
\quad
\b = \begin{bmatrix} 23 \\ 45 \end{bmatrix},
\quad
\c^\top = \begin{bmatrix} 2 & 1 & 0 & 0 \end{bmatrix},
\quad
B = \{x, y\},\ N = \{s_1, s_2\}.
\]
Kemudian,
\[
A_B = \begin{bmatrix} 3 & 1 \\ 5 & 3 \end{bmatrix},
\qquad
A_B^{-1} = \frac{1}{4} \begin{bmatrix} 3 & -1 \\ -5 & 3 \end{bmatrix}
= \begin{bmatrix} 0.75 & -0.25 \\ -1.25 & 0.75 \end{bmatrix}.
\]
Karena $A_N = I_2$ (kolom slack),
\[
A_B^{-1}\b = \frac14 \begin{bmatrix} 3(23) - 45 \\ -5(23) + 3(45) \end{bmatrix}
= \begin{bmatrix} 6 \\ 5 \end{bmatrix},
\qquad
A_B^{-1}A_N = A_B^{-1}
= \begin{bmatrix} 0.75 & -0.25 \\ -1.25 & 0.75 \end{bmatrix}.
\]
Bagian kendala dari kamus,
$\x_B = A_B^{-1}\b - A_B^{-1}A_N \x_N$, berbunyi
\[
\begin{aligned}
x &= 6 - 0.75\,s_1 + 0.25\,s_2, \\
y &= 5 + 1.25\,s_1 - 0.75\,s_2,
\end{aligned}
\]
tepat seperti baris yang ditampilkan. Untuk objektif,
$\c_B^\top A_B^{-1}\b = 2(6) + 1(5) = 17$, dan
\[
\c_B^\top A_B^{-1} = \begin{bmatrix} 2 & 1 \end{bmatrix}
\begin{bmatrix} 0.75 & -0.25 \\ -1.25 & 0.75 \end{bmatrix}
= \begin{bmatrix} 0.25 & 0.25 \end{bmatrix},
\qquad
\c_N^\top - \c_B^\top A_B^{-1}A_N
= \begin{bmatrix} -0.25 & -0.25 \end{bmatrix}.
\]
Dengan demikian, rumus tersebut memberikan
$\max\ 17 - 0.25 s_1 - 0.25 s_2$, yang mereproduksi seluruh kamus yang
ditampilkan. (Kedua biaya tereduksinya negatif, sehingga basis ini optimal.)

\exsol{8.7}{Biaya tereduksi variabel basis}{2}

Menerapkan rumus biaya tereduksi pada kolom-kolom basis berarti mengganti
$A_N$ dengan $A_B$:
\[
\c_B^\top - \c_B^\top A_B^{-1} A_B
= \c_B^\top - \c_B^\top I_m
= \c_B^\top - \c_B^\top = \mathbf{0}^\top,
\]
karena $A_B^{-1}A_B = I_m$ menurut definisi invers. Jadi, setiap variabel basis
memiliki biaya tereduksi tepat nol.

Interpretasi: kamus pada basis $B$ menyatakan $z$ dan setiap variabel basis
murni dalam variabel-variabel \emph{nonbasis}. Menyusun kamus berarti
menyubstitusikan
$\x_B = A_B^{-1}\b - A_B^{-1}A_N\x_N$ ke fungsi objektif, yang
mengeleminasi setiap variabel basis dari baris objektif; perhitungan di atas
merupakan pernyataan aljabar bahwa eliminasi ini eksak. Variabel basis tidak
dapat muncul pada baris objektif karena, jika muncul, baris tersebut belum
diselesaikan terhadap variabel basis dan sifat ``kamus'' (variabel basis dan
$z$ sebagai fungsi yang hanya bergantung pada variabel nonbasis) akan gagal.
Dengan kata lain, nilai variabel basis ditentukan oleh variabel nonbasis,
sehingga variabel basis tidak membawa laju perubahan $z$ yang independen.

\exsol{8.8}{Tanda biaya tereduksi dan aturan masuk}{2}

\textbf{(a)} Dengan menuliskan kembali objektif dari sudut pandang basis $B$,
\[
z = \c_B^\top A_B^{-1}\b
+ \left(\c_N^\top - \c_B^\top A_B^{-1} A_N\right)\x_N,
\]
diperoleh tepat baris objektif kamus: suatu konstanta ditambah fungsi linier
dari variabel nonbasis. Koefisien satu variabel nonbasis $x_j$ pada baris ini
adalah komponen ke-$j$ dari vektor biaya tereduksi, yaitu
$c_j - \c_B^\top A_B^{-1} A_{\{j\}}$. Jadi, ``pilih variabel nonbasis dengan
koefisien baris objektif positif'' dan ``pilih variabel nonbasis dengan biaya
tereduksi positif'' benar-benar merupakan aturan yang sama. Syarat berhentinya
(tidak ada koefisien positif; semua biaya tereduksi $\leq 0$) juga bertepatan.

\textbf{(b)} Menaikkan variabel nonbasis $x_j$ dari nol, sambil
mempertahankan variabel nonbasis lainnya pada nol, menggerakkan solusi sepanjang
sinar
\[
\x(t) = \bar\x + t\,\mathbf{d},
\qquad
\mathbf{d}_B = -A_B^{-1}A_{\{j\}},\quad d_j = 1,\quad d_k = 0
\text{ selain itu}.
\]
Secara geometris, ini adalah gerakan menjauh dari titik sudut saat ini sepanjang
sebuah sisi daerah layak: satu kendala aktif ($x_j = 0$) dilepaskan, sementara
kendala-kendala kesamaan tetap terpenuhi. Objektif berubah dengan laju yang sama
dengan biaya tereduksi, sehingga biaya tereduksi positif berarti sisi tersebut
menanjak. Uji rasio menghentikan gerakan di titik sudut tetangga, tempat suatu
variabel basis mencapai nol.

\textbf{(c)} Tableau menyimpan objektif sebagai persamaan
$z - \c^\top\x = 0$, sehingga entri yang dicatat di bawah variabel nonbasis
adalah \emph{negatif} dari koefisien kamusnya, yaitu negatif dari biaya
tereduksinya. Karena itu, mencari entri baris $z$ negatif dalam tableau sama
dengan mencari biaya tereduksi positif, dan ``berhenti ketika semua entri baris
$z$ tak negatif'' sama dengan ``berhenti ketika semua biaya tereduksi
$\leq 0$.'' Ketiga formulasi tersebut menyatakan satu pengujian dengan dua
konvensi tanda; tidak ada kontradiksi.

\exsol{8.9}{Biaya tereduksi yang negatif secara ketat menyiratkan optimum unik}{3}

(Solusi pilihan dalam buku, telah diverifikasi.) Tuliskan
$\bar{\c}_N^\top = \c_N^\top - \c_B^\top A_B^{-1} A_N$ dan
$z^* = \c^\top \bar\x = \c_B^\top A_B^{-1} \b$. Setiap $\x$ yang layak
memenuhi $A_B \x_B + A_N \x_N = \b$, jadi
$\x_B = A_B^{-1}\b - A_B^{-1}A_N \x_N$ dan, tepat seperti dalam teorema
optimalitas,
\[
\c^\top \x = z^* + \bar{\c}_N^\top \x_N .
\]
Ambil sembarang $\x \neq \bar\x$ yang layak.

\emph{Kasus 1: $\x_N \neq \mathbf{0}$.} Maka ada komponen $x_j > 0$ dengan
$j \in N$, dan karena setiap entri dari $\bar{\c}_N$ benar-benar negatif,
$\bar{\c}_N^\top \x_N < 0$. Karena itu, $\c^\top \x < z^*$, sehingga $\x$
tidak optimal.

\emph{Kasus 2: $\x_N = \mathbf{0}$.} Maka $A_B \x_B = \b$, dan karena $A_B$
invertibel, ini memaksa $\x_B = A_B^{-1}\b$, yakni $\x = \bar\x$, yang
bertentangan dengan $\x \neq \bar\x$.

Oleh karena itu, setiap titik layak selain $\bar\x$ memiliki nilai objektif
yang secara ketat lebih kecil daripada $z^*$: solusi $\bar\x$ optimal dan tidak
ada solusi optimal lainnya.
